\documentclass[11pt,a4paper]{article}

\usepackage[a4paper,margin=0.85in]{geometry}
\usepackage[T1]{fontenc}
\usepackage{lmodern}
\usepackage{microtype}
\usepackage{graphicx}
\usepackage[table]{xcolor}
\usepackage{tabularx}
\usepackage{array}
\usepackage{fancyhdr}
\usepackage[hidelinks]{hyperref}

\definecolor{saidlblue}{HTML}{163A63}
\definecolor{saidlgrey}{HTML}{4D5966}
\definecolor{rulegrey}{HTML}{C9D0D7}

\setlength{\parindent}{0pt}
\setlength{\parskip}{0.4\baselineskip}
\pagestyle{fancy}
\fancyhf{}
\renewcommand{\headrulewidth}{0pt}
\fancyfoot[C]{\footnotesize\color{saidlgrey}SAiDL CTE Course Handout \textbullet\ First Semester 2026 to 2027}

\newcommand{\infoblock}[2]{%
  \textbf{\color{saidlblue}#1:} #2\par
}

\newcolumntype{L}{>{\centering\arraybackslash\bfseries\color{saidlblue}}p{1.55cm}}
\newcolumntype{Y}{>{\raggedright\arraybackslash}X}

\begin{document}

\begin{minipage}[c]{0.14\textwidth}
  \centering
  \includegraphics[width=1.65cm]{saidl-logo.jpeg}
\end{minipage}%
\begin{minipage}[c]{0.72\textwidth}
  \centering
  {\large\bfseries Birla Institute of Technology and Science, Pilani}\par
  {\normalsize\bfseries K. K. Birla Goa Campus}\par
  {\normalsize\bfseries Center for Technical Education}\par
  \vspace{0.15em}
  {\Large\bfseries\color{saidlblue}SAiDL}\par
  {\small\bfseries First Semester, 2026 to 2027}
\end{minipage}%
\begin{minipage}[c]{0.14\textwidth}\mbox{}\end{minipage}

\vspace{0.35em}
{\color{rulegrey}\rule{\textwidth}{0.8pt}}

\begin{center}
  {\LARGE\bfseries Course Title: This ML Course Is All You Need}
\end{center}

\vspace{-0.3em}
{\small
\infoblock{Instructors}{To be announced}
\infoblock{Mentors}{To be announced}
\infoblock{Course Description}{An introduction to machine learning from regression and neural-network fundamentals through computer vision, attention, transformers, generative AI, and reinforcement learning.}
\infoblock{Prerequisites}{Basic proficiency in Python; no prior machine-learning experience required.}
}

\vspace{0.25em}
\begin{center}
  {\Large\bfseries\color{saidlblue}Course Structure}
\end{center}
\vspace{-0.35em}

{\small
\renewcommand{\arraystretch}{1.08}
\setlength{\tabcolsep}{6pt}
\begin{tabularx}{\textwidth}{|L|Y|}
  \hline
  \rowcolor{saidlblue}
  \multicolumn{1}{|c|}{\color{white}\bfseries Lecture} &
  \multicolumn{1}{c|}{\color{white}\bfseries Agenda} \\
  \hline
  L1 & \textbf{Overview:} What is ML?; real-life applications; linear and logistic regression. \\
  \hline
  L2 & \textbf{How ML Models Learn:} Supervised learning; loss functions; gradient descent; classifiers. \\
  \hline
  L3 & \textbf{Building a Neural Network:} Neural networks; backpropagation; activation functions. \\
  \hline
  L4 & \textbf{Making Neural Networks Learn:} Optimization; dropout; initialization; regularization. \\
  \hline
  L5 & \textbf{Computer Vision:} Convolutional neural networks; pooling; feature maps. \\
  \hline
  L6 & \textbf{Sequential Models:} ResNets; recurrent neural networks; long short-term memory networks. \\
  \hline
  L7 & \textbf{Attention:} Attention; queries, keys, and values; self-attention; cross-attention; tokenization. \\
  \hline
  L8 & \textbf{Transformer Architecture:} Multi-head attention; positional encoding; the Transformer. \\
  \hline
  L9 & \textbf{ML Without Labels:} Unsupervised learning; K-means clustering; pre-training; self-supervised learning. \\
  \hline
  L10 & \textbf{Generative Models:} Generative modeling; GPTs; variational autoencoders; Vision Transformers. \\
  \hline
  L11 & \textbf{Modern AI Systems:} RLHF; prompt engineering; ReAct; chain-of-thought reasoning. \\
  \hline
  L12 & \textbf{Reinforcement Learning:} Agents; Markov decision processes; value functions; Q-learning. \\
  \hline
\end{tabularx}

\vspace{0.35em}
\infoblock{Contact Information}{To be announced}
}

\end{document}
